%=====================================================================
\chapter{Data: Types, Sources and Quality}\label{ch:data}
%=====================================================================
\locator{1}{Part I --- Foundations of Data}

\begin{outcomes}
By the end of this chapter you will be able to:
\begin{tight}
  \item \textbf{Classify} data as structured, semi-structured or unstructured, and state what each implies for storage and analysis.
  \item \textbf{Assign} a variable to its correct measurement scale --- nominal, ordinal, interval or ratio --- and explain which operations that permits.
  \item \textbf{Evaluate} a dataset against the six dimensions of data quality.
  \item \textbf{Distinguish} the three mechanisms of missingness and explain why the distinction changes what you may do about it.
  \item \textbf{Describe} the characteristic data of each of the four running domains.
\end{tight}
\end{outcomes}

\begin{prereq}
\chref{lifecycle} --- the framing questions, especially ``one row is one what?'',
which this chapter makes precise.
\end{prereq}

\begin{keyterms}
observation $\bullet$ feature $\bullet$ structured / semi-structured / unstructured
$\bullet$ nominal, ordinal, interval, ratio $\bullet$ tidy data $\bullet$ metadata
$\bullet$ provenance $\bullet$ data quality dimensions $\bullet$ MCAR, MAR, MNAR
$\bullet$ granularity
\end{keyterms}

\section{The Shape of a Dataset}

\begin{definitionbox}[Observation and Feature]
An \term{observation} (row, record, example, instance) is one thing you measured.
A \term{feature} (column, variable, attribute, field) is one property measured
about every thing. The number of observations is written $n$; the number of
features $p$.
\end{definitionbox}

The apparently trivial question ``what is one row?'' is the source of more
confusion than any other in applied work, because the same underlying reality
admits several answers.

\begin{worked}[The Same Data, Four Granularities]
A university has swipe-card records for lecture attendance.

\begin{tight}
  \item \emph{One row per swipe:} $n \approx 400{,}000$. Features: student, room,
        timestamp. This is the raw \term{granularity}.
  \item \emph{One row per student per lecture:} $n \approx 120{,}000$. Feature:
        attended yes/no.
  \item \emph{One row per student per week:} $n \approx 24{,}000$. Feature:
        lectures attended out of 5.
  \item \emph{One row per student:} $n \approx 2{,}000$. Feature: overall
        attendance percentage.
\end{tight}

\textbf{Why it matters.} If the frame says ``predict which \emph{student} fails'',
the modelling table must have one row per student --- so the other three must be
aggregated up. Every aggregation destroys information (row 4 cannot tell you
\emph{when} attendance collapsed) and every disaggregation invents none. Choose
the coarsest granularity that still contains the signal, and do the aggregation
deliberately rather than by accident.
\end{worked}

\section{Three Kinds of Data}

\dsfig{%
\begin{tikzpicture}[
  b/.style={rounded corners=4pt, draw=#1, line width=1.2pt, fill=#1!7,
            text width=4.4cm, align=center, font=\scriptsize, inner sep=6pt},
  ttl/.style={font=\bfseries\small}]

  \node[b=primary] (s) at (0,0) {%
    {\color{primary}\bfseries\small \faIcon{table}~STRUCTURED}\\[3pt]
    Fixed schema, rows and columns.\\[2pt]
    \emph{Tables, CSV, SQL databases, spreadsheets}\\[3pt]
    {\color{primary!60!black}Ready for analysis; roughly 20\% of the world's data}};

  \node[b=secondary] (m) at (5.3,0) {%
    {\color{secondary}\bfseries\small \faIcon{code}~SEMI-STRUCTURED}\\[3pt]
    Tags or keys, but no fixed schema.\\[2pt]
    \emph{JSON, XML, logs, sensor payloads}\\[3pt]
    {\color{secondary!60!black}Must be flattened before most analysis}};

  \node[b=tealhl] (u) at (10.6,0) {%
    {\color{tealhl}\bfseries\small \faIcon{file-image}~UNSTRUCTURED}\\[3pt]
    No predefined structure at all.\\[2pt]
    \emph{Text, images, audio, video, PDFs}\\[3pt]
    {\color{tealhl!60!black}Needs deep learning (Part IV) to become features}};

  \draw[-{Stealth[length=2.5mm]}, gray!45, line width=6pt, opacity=0.35]
        (-2.4,-1.75) -- (13.0,-1.75);
  \node[font=\scriptsize\itshape, text=gray!60!black] at (5.3,-2.15)
       {increasing volume ~~$\longrightarrow$~~ increasing effort to make usable};
\end{tikzpicture}}%
{The three kinds of data. Most organisational data is not the tidy table beginners picture;
the majority is semi-structured logs and unstructured documents, which is why Chapters 7,
20 and 23 exist.}{data-kinds}

\section{Measurement Scales}

Knowing that a column contains numbers is not enough. Postcodes are numbers you
must never average; temperatures in Celsius are numbers you may average but must
not double.

\begin{definitionbox}[The Four Measurement Scales]
\term{Nominal} --- named categories, no order (blood group, protocol, department).\\
\term{Ordinal} --- ordered categories, unequal or unknown gaps (grade A/B/C, severity low/medium/high).\\
\term{Interval} --- numeric, equal gaps, \emph{arbitrary zero} (Celsius, calendar year).\\
\term{Ratio} --- numeric, equal gaps, \emph{true zero} (count, duration, bytes, kelvin).
\end{definitionbox}

\rowcolors{2}{white}{defbg}
\begin{longtable}{@{}L{2.0cm}L{2.6cm}L{2.3cm}L{3.0cm}L{4.2cm}@{}}
\toprule
\rowcolor{primary!15}
\textbf{Scale} & \textbf{Legal operations} & \textbf{Centre} & \textbf{Example} & \textbf{Nonsense you must avoid} \\
\midrule
\endhead
Nominal & $=$, $\neq$ & Mode & TCP / UDP / ICMP & ``the average protocol is 1.7'' \\
Ordinal & $=$, $<$, $>$ & Median & Low / Med / High risk & ``High is exactly twice Medium'' \\
Interval & $+$, $-$ & Mean & Temperature in $^\circ$C & ``$20^\circ$C is twice as hot as $10^\circ$C'' \\
Ratio & $+$, $-$, $\times$, $\div$ & Mean & Packets per second & (none --- all operations valid) \\
\bottomrule
\end{longtable}

\begin{alertbox}[The Encoded-Category Trap]
The moment you encode \texttt{\{TCP, UDP, ICMP\}} as \texttt{\{1, 2, 3\}} so a
model will accept it, every algorithm downstream will happily compute that
$\text{UDP} - \text{TCP} = \text{TCP}$ and that ICMP is ``larger'' than the others.
This is why \chref{wrangling} teaches one-hot encoding: it preserves the nominal
scale instead of silently promoting it to ratio.
\end{alertbox}

\section{Tidy Data}

\begin{definitionbox}[Tidy Data]
A dataset is \term{tidy} when (1) each variable forms a column, (2) each
observation forms a row, and (3) each type of observational unit forms a table.
Almost every analysis tool assumes tidy input, and almost no real dataset arrives
that way.
\end{definitionbox}

\dsfig{%
\begin{tikzpicture}[font=\scriptsize,
  hdr/.style={fill=accent!18, draw=accent!60, text=accent!50!black, font=\scriptsize\bfseries,
              minimum width=1.55cm, minimum height=0.48cm},
  cel/.style={fill=white, draw=gray!45, minimum width=1.55cm, minimum height=0.48cm},
  ghdr/.style={fill=greenhl!18, draw=greenhl!60, text=greenhl!45!black, font=\scriptsize\bfseries,
              minimum width=1.35cm, minimum height=0.48cm},
  gcel/.style={fill=white, draw=gray!45, minimum width=1.35cm, minimum height=0.48cm}]

% ---- messy (wide) ----
\node[font=\bfseries\small, text=accent] at (2.3,1.35) {Messy: weeks hidden in column names};
\foreach \c/\t in {0/student, 1/wk1, 2/wk2, 3/wk3}
  \node[hdr] at (\c*1.55,0.65) {\t};
\foreach \r/\a/\b/\c/\d in {0/S101/5/4/2, 1/S102/3/0/0}{
  \pgfmathsetmacro{\yy}{0.65-(\r+1)*0.48}
  \node[cel] at (0,\yy) {\a}; \node[cel] at (1.55,\yy) {\b};
  \node[cel] at (3.10,\yy) {\c}; \node[cel] at (4.65,\yy) {\d};}

\draw[-{Stealth[length=3mm]}, primary!60, line width=2pt] (5.6,-0.1) -- (7.0,-0.1);
\node[font=\scriptsize\itshape, text=primary!70!black] at (6.3,0.35) {tidy};

% ---- tidy (long) ----
\node[font=\bfseries\small, text=greenhl!50!black] at (9.6,1.35) {Tidy: week is a variable};
\foreach \c/\t in {0/student, 1/week, 2/attended}
  \node[ghdr] at (7.9+\c*1.35,0.65) {\t};
\foreach \r/\a/\b/\c in {0/S101/1/5, 1/S101/2/4, 2/S101/3/2, 3/S102/1/3, 4/S102/2/0, 5/S102/3/0}{
  \pgfmathsetmacro{\yy}{0.65-(\r+1)*0.48}
  \node[gcel] at (7.9,\yy) {\a}; \node[gcel] at (9.25,\yy) {\b}; \node[gcel] at (10.6,\yy) {\c};}
\end{tikzpicture}}%
{Messy versus tidy layout of the same attendance data. On the left, \emph{week} is not a
variable but three column names, so you cannot filter, group or plot by it. Reshaping to
the right is the single most common preparation step in real projects.}{tidy-data}

\section{Data Quality}

\begin{definitionbox}[The Six Dimensions of Data Quality]
\term{Completeness} --- is anything missing?
\term{Accuracy} --- does it match reality?
\term{Consistency} --- do different sources agree?
\term{Timeliness} --- is it current enough for the decision?
\term{Validity} --- does it satisfy the rules of its own schema?
\term{Uniqueness} --- is each real-world entity represented once?
\end{definitionbox}

\begin{worked}[Auditing a Student Records Extract]
An extract of 2{,}000 student records is handed over. Applying the six dimensions
in order:

\smallskip
\textbf{Completeness:} \texttt{prior\_qualification} is empty for 640 rows (32\%).
Not fatal, but no model may treat it as reliably present.

\textbf{Accuracy:} 12 students have \texttt{age = 0}. Compare against date of
birth: these are placeholder values, not newborns.

\textbf{Consistency:} the registry lists 2{,}000 students; the VLE lists 2{,}043.
The extra 43 are staff test accounts. Left in, they would be modelled as students
who never fail.

\textbf{Timeliness:} \texttt{attendance\_pct} was computed in week 12 but the
model is to run in week 4. This is a \emph{framing} violation (\chref{lifecycle},
question 4), discovered here.

\textbf{Validity:} 3 rows carry \texttt{final\_grade = "Z"}, which is not in the
permitted grade set.

\textbf{Uniqueness:} 8 students appear twice, having transferred between
programmes. Deduplicating naively by student ID would silently drop one of their
two enrolments.

\smallskip
\textbf{Verdict:} the dataset is usable, but only after documented decisions on
each of the six. The audit takes an afternoon; skipping it costs weeks.
\end{worked}

\subsection{Missing Data: Three Mechanisms}

Why a value is missing determines what you are allowed to do about it. This is
one of the few places in the course where a philosophical distinction has an
immediate practical consequence.

\rowcolors{2}{white}{lightbg}
\begin{longtable}{@{}L{2.6cm}L{5.4cm}L{6.4cm}@{}}
\toprule
\rowcolor{primary!15}
\textbf{Mechanism} & \textbf{Meaning} & \textbf{Example and consequence} \\
\midrule
\endhead
MCAR \newline \emph{completely at random} & Missingness is unrelated to anything, observed or not & A sensor drops packets when the network is congested, independent of the reading. Dropping those rows is safe, just wasteful. \\
MAR \newline \emph{at random} & Missingness depends on \emph{observed} variables & Part-time students skip the optional survey more often. Missingness is predictable from \texttt{enrolment\_type}, so imputation conditional on it is defensible. \\
MNAR \newline \emph{not at random} & Missingness depends on the \emph{unobserved value itself} & Students who are failing stop submitting the self-assessment. The missingness \emph{is} the signal. Imputing it destroys the very information you needed. \\
\bottomrule
\end{longtable}

\begin{conceptbox}[The MNAR Rule of Thumb]
If you suspect MNAR, do not impute --- \emph{encode}. Add a boolean column
\texttt{was\_missing} and let the model use the fact of absence as a feature. In
learning analytics and in security logging, absence is very often the strongest
signal in the dataset.
\end{conceptbox}

\section{Provenance and Metadata}

\begin{definitionbox}[Provenance]
The documented history of a dataset: where it came from, who produced it, when,
under what definitions, and what has been done to it since. A dataset without
provenance is an anecdote.
\end{definitionbox}

Minimum metadata to record for every dataset you use in this course:
\begin{tight}
  \item Source system and extraction date.
  \item The definition of each column, \emph{in the words of the people who
        generate it} --- ``active user'' means something specific and probably
        surprising.
  \item Known quality issues and the decisions you took about them.
  \item Licence and permitted use, especially for personal data
        (see \chref{ethics}).
\end{tight}

%---------------------------------------------------------------------
\begin{fourlenses}{Characteristic Data}
\lensLA{\emph{Structured:} enrolment records, grades, demographics.
\emph{Semi-structured:} VLE clickstream (JSON events with timestamps).
\emph{Unstructured:} forum posts, essay submissions. \emph{Key quality risk:}
uniqueness --- one human may hold several accounts across systems, and joining
them wrongly either fragments or fabricates a student. \emph{Typical missingness:}
MNAR, because disengaged students stop generating data at exactly the moment you
most need it.}

\lensPM{\emph{Structured:} issue trackers, timesheets, commit counts.
\emph{Semi-structured:} CI/CD pipeline logs. \emph{Unstructured:} retrospective
notes, requirement documents. \emph{Key quality risk:} accuracy --- effort is
self-reported and rounded to comfortable numbers. \emph{Typical missingness:} MAR,
since teams under pressure log less, and ``under pressure'' is itself observable
from the sprint data.}

\lensIOT{\emph{Structured:} time-stamped numeric telemetry, one row per sensor
per interval. \emph{Semi-structured:} MQTT payloads with device-specific fields.
\emph{Unstructured:} camera and audio streams. \emph{Key quality risk:} validity
and timeliness --- clock drift across a fleet makes timestamps disagree, and a
reading that arrives two hours late is not the same as a reading that is on time.
\emph{Typical missingness:} MCAR from packet loss, but MNAR when a failing sensor
stops reporting --- and the failing ones are the ones you care about.}

\lensSEC{\emph{Structured:} NetFlow records, authentication events.
\emph{Semi-structured:} syslog, firewall and SIEM logs. \emph{Unstructured:}
malware binaries, phishing email bodies. \emph{Key quality risk:} completeness and
consistency --- log sources fail silently, and an attacker's first act is often to
clear or disable logging. \emph{Typical missingness:} MNAR by design, because the
adversary is deliberately creating the gap.}
\end{fourlenses}

%---------------------------------------------------------------------
\begin{lab}[A Ten-Minute Data Audit]
Run this on any new dataset before doing anything else.

\begin{lstlisting}[language=Python]
import pandas as pd

def audit(df):
    print(f"shape: {df.shape[0]} rows x {df.shape[1]} columns\n")

    print("-- completeness (% missing) --")
    print((df.isna().mean() * 100).round(1).sort_values(ascending=False), "\n")

    print("-- uniqueness --")
    print(f"duplicate rows: {df.duplicated().sum()}\n")

    print("-- validity: candidate placeholder values --")
    for c in df.select_dtypes("number"):
        odd = df[c].isin([0, -1, -999]).sum()
        if odd:
            print(f"  {c}: {odd} rows equal 0/-1/-999")

    print("\n-- scale check: how many distinct values? --")
    print(df.nunique().sort_values())   # few distinct + numeric => probably nominal

audit(pd.read_csv("your_data.csv"))
\end{lstlisting}

\textbf{Read the last output carefully.} A numeric column with only four distinct
values is almost certainly a category wearing a number's clothing.
\end{lab}

%---------------------------------------------------------------------
\begin{checkpoint}
\begin{tightnum}
  \item Classify each by measurement scale: (a)~student ID, (b)~satisfaction on a
        1--5 Likert scale, (c)~CPU temperature in $^\circ$C, (d)~number of failed
        logins.
  \item A sensor reports \texttt{-999} when it cannot obtain a reading. What are
        the two things that will go wrong if you leave those values in and compute
        a mean? What is the missingness mechanism most likely to be?
  \item Attendance data is stored one row per swipe. Your frame requires one row
        per student per week. Is that aggregation or disaggregation, and what
        information is lost?
\end{tightnum}
\end{checkpoint}

%---------------------------------------------------------------------
\begin{chaptersummary}
\begin{tight}
  \item An observation is a row, a feature is a column, and deciding what one row
        represents (granularity) is a framing decision, not a technical one.
  \item Data is structured, semi-structured or unstructured; most real data is not
        the first.
  \item The measurement scale --- nominal, ordinal, interval, ratio --- determines
        which operations are meaningful. Encoding a category as a number does not
        make arithmetic on it valid.
  \item Tidy data means one variable per column, one observation per row. Most
        data arrives untidy.
  \item Quality has six dimensions; audit all six before modelling.
  \item Missingness is MCAR, MAR or MNAR. Under MNAR, encode the absence rather
        than imputing it --- the gap is the signal.
\end{tight}
\end{chaptersummary}

\begin{reviewq}
  \item \textbf{(MCQ)} Which is a ratio-scale variable? (a)~Calendar year
        (b)~Temperature in $^\circ$C (c)~Number of packets sent (d)~Severity rated
        low/medium/high.
  \item \textbf{(MCQ)} JSON log events are best described as: (a)~structured
        (b)~semi-structured (c)~unstructured (d)~tidy.
  \item \textbf{(Short)} Explain why it is invalid to say ``$20^\circ$C is twice as
        warm as $10^\circ$C'' but valid to say ``200 packets is twice as many as
        100''.
  \item \textbf{(Short)} Define MNAR and give an example from cybersecurity in
        which the missingness itself is the most informative feature available.
  \item \textbf{(Short)} A dataset has no provenance record. List two specific
        risks this creates for a project that reaches deployment.
  \item \textbf{(Applied)} You are given swipe-card, VLE-clickstream and grade data
        for 2{,}000 students. Design the granularity of your modelling table for the
        frame ``predict in week 4 whether a student will fail'', and list the
        aggregations required to get there.
  \item \textbf{(Applied)} For an IoT fleet of 400 motors reporting every 10
        seconds, audit the dataset against all six quality dimensions: for each,
        state one concrete check you would run.
\end{reviewq}
